%%%% RESUMO
%%
%% Apresentação concisa dos pontos relevantes de um texto, fornecendo uma visão rápida e clara do conteúdo e das conclusões do
%% trabalho.

\begin{resumoutfpr}
O desenvolvimento colaborativo de \textit{software}, especialmente na indústria de jogos digitais, envolve o uso intensivo de arquivos binários de grande porte, como modelos tridimensionais, texturas e áudios. Os sistemas tradicionais de controle de versão, ferramentas que registram e gerenciam o histórico de alterações em arquivos ao longo do tempo, não lidam de forma eficiente com esse tipo de arquivo, pois armazenam cópias completas a cada modificação, mesmo quando apenas pequenas partes do conteúdo foram alteradas. Soluções existentes no mercado atenuam parcialmente esse problema, mas ainda tratam os arquivos binários como blocos indivisíveis, sem analisar seu conteúdo interno. Este trabalho propõe uma abordagem de versionamento que armazena apenas as partes efetivamente modificadas de um arquivo binário tridimensional, aproveitando sua estrutura interna para identificar e isolar cada segmento de dados. Para detectar alterações, aplica-se uma função de SHA-256 sobre cada segmento, gerando uma assinatura digital única que permite comparar versões com precisão. O armazenamento diferencial é realizado por meio de um sistema de arquivos moderno, como o Btrfs que, ao registrar uma modificação, preserva os dados anteriores sem sobrescrevê-los, compartilhando os blocos inalterados entre as versões. Os experimentos compararam o método proposto com a abordagem convencional de armazenamento de arquivos grandes em repositórios de código, em cenários de modificação progressiva de um modelo tridimensional. Os resultados iniciais indicam reduções expressivas no espaço consumido, chegando a economias de até 98\% por versão nos cenários com poucas alterações, e de aproximadamente 39\% no espaço total acumulado ao longo de múltiplas versões.
\end{resumoutfpr}
